\section{Постановка задачи}\label{sec:problem}

\subsection{Модель и функция потерь}

Рассматривается однослойная полносвязная сеть без смещений
\begin{equation}
  f(x; W_0, W_1) = W_1\,\sigma(W_0 x),
\end{equation}
где $x\in\mathbb{R}^{d_{in}}$, $W_0\in\mathbb{R}^{h\times d_{in}}$,
$W_1\in\mathbb{R}^{d_{out}\times h}$, а $\sigma$ --- поэлементная функция
активации (tanh, ReLU либо тождественная). Всюду далее
$d_{in}=d_{out}=h=16$, что даёт $2\cdot 16\cdot 16 = 512$ обучаемых параметров.

Функция потерь --- среднеквадратичная ошибка на обучающей выборке
$\mathcal{D}=\{(x_i,y_i)\}_{i=1}^{N}$:
\begin{equation}
  L(W_0,W_1) = \frac{1}{N d_{out}}\sum_{i=1}^{N}\bigl\|f(x_i;W_0,W_1)-y_i\bigr\|^2 .
\end{equation}

\subsection{Семейство оптимизаторов}

Пусть $g_t^{(\ell)} = \nabla_{W_\ell} L\bigl(W^{(t)}\bigr) + \lambda_{wd} W^{(t)}_\ell$ ---
градиент по весам слоя $\ell\in\{0,1\}$ с учётом $L_2$-регуляризации
(затухание весов реализовано добавлением к градиенту, не в развязанной форме
AdamW). Рассматриваются четыре правила обновления.

\smallskip
\noindent\textbf{SGD:}\quad $W^{(t+1)}_\ell = W^{(t)}_\ell - \eta_\ell g^{(\ell)}_t$.

\smallskip
\noindent\textbf{SGD с инерцией (heavy ball):}
\begin{equation}\label{eq:hb}
  v^{(\ell)}_{t+1} = \beta v^{(\ell)}_t + g^{(\ell)}_t,
  \qquad
  W^{(t+1)}_\ell = W^{(t)}_\ell - \eta_\ell v^{(\ell)}_{t+1}.
\end{equation}

\smallskip
\noindent\textbf{RMSprop:}\quad
$s_{t+1} = \beta_2 s_t + (1-\beta_2) g_t^{\odot 2}$,\quad
$W^{(t+1)} = W^{(t)} - \eta\, g_t / (\sqrt{s_{t+1}}+\varepsilon)$.

\smallskip
\noindent\textbf{Adam:}\quad
$m_{t+1} = \beta_1 m_t + (1-\beta_1) g_t$,\quad
$s_{t+1} = \beta_2 s_t + (1-\beta_2) g_t^{\odot 2}$, с коррекцией смещения
$\hat m = m_{t+1}/(1-\beta_1^{t+1})$, $\hat s = s_{t+1}/(1-\beta_2^{t+1})$ и
обновлением $W^{(t+1)} = W^{(t)} - \eta\,\hat m/(\sqrt{\hat s}+\varepsilon)$,
$\varepsilon = 10^{-8}$.

Все операции выполняются послойно, что позволяет задавать отдельные шаги
обучения $\eta_0$ и $\eta_1$ для входного и выходного слоёв.

\subsection{Классификация режимов обучения}

Пусть $\tilde L(t) = L(W^{(t)})/L(W^{(0)})$ --- нормализованный loss. В отличие
от предыдущей НИР, где использовалось бинарное деление на сходимость и
расходимость, здесь вводится классификация на четыре режима, поскольку при
наличии инерции и адаптивных методов ограниченные незатухающие колебания
занимают значительную долю пространства гиперпараметров.

\begin{definition}[Режимы]\label{def:regimes}
Пусть $T$ --- число итераций, $\mathcal{T}=\{T-64,\dots,T-1\}$ --- хвост
траектории. Обозначим
$m_{\text{кон}} = \frac{1}{20}\sum_{t=T-20}^{T-1}\tilde L(t)$,
$m_{\text{нач}} = \frac{1}{20}\sum_{t\in\mathcal{T},\,t<T-44}\tilde L(t)$,
а через $\mathrm{rel} = \operatorname{std}_{\mathcal{T}}\tilde L / \operatorname{mean}_{\mathcal{T}}\tilde L$ ---
относительный разброс на хвосте. Тогда:
\begin{itemize}[leftmargin=*,topsep=2pt,itemsep=1pt]
  \item \textbf{NaN/Inf}, если хотя бы на одной итерации $\tilde L(t)$ не является
        конечным числом;
  \item \textbf{расходимость}, если $\max_t \tilde L(t) > 10^{6}$, но переполнения
        не произошло;
  \item \textbf{сходимость}, если $m_{\text{кон}} < 1$ и при этом либо
        $\mathrm{rel} < 10^{-3}$ (траектория вышла на стационар), либо
        $m_{\text{кон}} < 0.99\, m_{\text{нач}}$ (loss продолжает убывать);
  \item \textbf{осцилляции/периодичность} --- во всех остальных случаях:
        траектория ограничена, но не сходится (циклы периода 2 и выше,
        хаотические колебания, стационарные плато выше начального уровня).
\end{itemize}
\end{definition}

Условие $m_{\text{кон}}<1$ совпадает с критерием сходимости предыдущей НИР;
дополнительные условия отделяют от сходимости ограниченные незатухающие режимы.
Границей обучаемости далее называется граница множества
$\{\text{сходимость}\}$ в пространстве гиперпараметров.

\subsection{Сканирование и измерение размерности}

Для пары гиперпараметров $(p,q)$ строится прямоугольная сетка $n\times n$
(логарифмическая по шагу обучения и по затуханию весов, равномерная по
$\beta$, $\beta_1$, $\beta_2$ и по размеру батча). Для каждого узла выполняется
$T$ итераций обучения из одной и той же фиксированной инициализации, после чего
узел относится к одному из четырёх режимов согласно
определению~\ref{def:regimes}. Результат --- матрица режимов $R\in\{0,1,2,3\}^{n\times n}$.

Граница извлекается морфологически. Пусть $B_{ij}=[\,R_{ij}=\text{сходимость}\,]$.
Тогда
\begin{equation}
  \partial B = \mathrm{Dilation}(B)\ \oplus\ \mathrm{Erosion}(B),
\end{equation}
где $\oplus$ --- симметрическая разность (XOR), после чего применяется
скелетонизация до толщины в один пиксель.

Box-counting размерность оценивается стандартно~\cite{falconer2014}: для
последовательности размеров ячеек $\varepsilon_k = 2^k$, $k=1,\dots,\lfloor\log_2 n\rfloor-2$,
подсчитывается число $N(\varepsilon_k)$ непустых ячеек, покрывающих $\partial B$, и
\begin{equation}\label{eq:box}
  D_{\mathrm{box}} = -\,\mathrm{slope}\bigl(\text{МНК-регрессия } \log N(\varepsilon)\ \text{по}\ \log\varepsilon\bigr).
\end{equation}
В качестве погрешности $\sigma_D$ приводится стандартная ошибка наклона
регрессии; она характеризует только качество линейной аппроксимации в
координатах $\log\varepsilon$--$\log N$ и не учитывает систематическую
зависимость оценки от разрешения сетки, которая обсуждается отдельно
в разделе~\ref{sec:limitations}.

\subsection{Что именно проверяется}

Формально исследуются следующие гипотезы:
\begin{enumerate}[leftmargin=*,topsep=2pt,itemsep=1pt]
  \item фрактальность границы сохраняется в сечении $(\eta,\beta)$, то есть
        добавление инерции не делает границу гладкой кривой;
  \item затухание весов и увеличение размера батча уменьшают $D_{\mathrm{box}}$;
  \item адаптивные методы (Adam, RMSprop) дают меньшую $D_{\mathrm{box}}$, чем SGD,
        в сопоставимых сечениях;
  \item структурированные данные не разрушают фрактальность границы.
\end{enumerate}
